% document formatting
\documentclass[10pt]{article}
% \usepackage[UTF8]{inputenc}
\usepackage[UTF8, scheme=plain]{ctex}
\usepackage[left=1in,right=1in,top=1in,bottom=1in]{geometry}
\usepackage[T1]{fontenc}
\usepackage{xcolor}

% math symbols, etc.
\usepackage{amsmath, amsfonts, amssymb, amsthm}

% lists
\usepackage{enumerate}
\usepackage[table,xcdraw]{xcolor}

% images
\usepackage{graphicx} % for images
\usepackage{tikz}

% code blocks
% \usepackage{minted, listings} 

% verbatim greek
\usepackage{alphabeta}

\graphicspath{{./assets/images/Week 3}}

\title{02-741 Week 3 \\ \large{Generative AI for Biomedicine}}
 
\author{Aidan Jan}

\date{\today}

\begin{document}
\maketitle

\section*{Diffusion Models}
Diffusion models are state-of-the-art models for (continuous) data generation
\begin{itemize}
	\item Time series data
	\item Audio / speech
	\item 3D objects (coordinates), e.g., Molecule structures
\end{itemize}

\subsection*{Data Representation}
\begin{itemize}
	\item Images: $x \in \mathbb{R}^{w \times h}$
	\item Speech: $x \in \mathbb{R}^t$
	\item Molecular structure: $x \in \mathbb{R}^{N \times 3}$
\end{itemize}
In each case, we want a probabilistic model $p(x)$ to fit the data and generate samples

\subsection*{Probabilistic Generative Process}
\begin{itemize}
	\item Start from initial distribution $p_T(x_T)$ (e.g., Gaussian $N(0, I)$).  This is just noise.
	\item Generate sligntly improved data: $x_{t - 1} \sim p_{t - 1}(x_{t - 1} | x_t)$
	\item Final data: $x_0 \sim p_0 (x_0 | x_1)$
\end{itemize}
The generation process is a Markov chain.
\begin{center} 
	\includegraphics*[width=0.9\textwidth]{W3_1.png} 
\end{center}
Learning the generative model is equivalent to learning the parameters for each $p_{t - 1}(x_{t - 1} | x_t)$.
\begin{itemize}
	\item It is difficult to directly construct a series of probability distributions
\end{itemize}

\section*{Diffusion Model}
\begin{enumerate}
	\item Forward Noising Process
	\begin{itemize}
	    \item Synthesize data for training by repeatedly adding noise each step
    \end{itemize}
	\item Reverse Generation Process
	\begin{itemize}
	    \item Start from complete noise, generate clean data
	    \item Also known as denoising
    \end{itemize}
\end{enumerate}

\subsection*{(Forward) Noising Process}
\begin{itemize}
	\item Adding Gaussian noise at each time step $t$
	\[x_t \sim q(x_t | x_{t - 1}) = N(\sqrt{\alpha_t} x_{t - 1}, \beta_t I), \alpha_t = 1 - \beta_t\]
    equivalently,
    \[x_t = \sqrt{\alpha_t} x_{t - 1} + \sqrt{\beta_t} \epsilon, \epsilon \sim N(0, I)\]
\end{itemize}
\begin{center} 
	\includegraphics*[width=0.9\textwidth]{W3_2.png} 
\end{center}
Consequence: with sufficient large $T$, the result will be Gaussian noise.

\subsection*{(Reverse) Generation / Denoising Process}
\begin{center} 
	\includegraphics*[width=0.9\textwidth]{W3_3.png} 
\end{center}
Learning problem: how to find parameters of $p$ to match the sequence of noised images?

\subsubsection*{Graphical Model for Diffusion}
\begin{center} 
	\includegraphics*[width=0.9\textwidth]{W3_4.png} 
\end{center}

\subsection*{Denoising Diffusion Probabilistic Models (DDPM)}
\begin{itemize}
	\item Reverse Process (Generation), aka. denoising
	\begin{align*}
        p_\theta &= \prod_{t = 1}^T p_\theta (x_{t - 1} | x_t) \\
        p_\theta(x_{t - 1}|x_t) &= N(x_{t - 1}; \mu_\theta (x_t, t), \Sigma_\theta(x_t, t))
    \end{align*}
    the $\mu_\theta$ and $\Sigma_\theta$ functions are predicted by neural network
    \item Forward Process (Diffusion), aka noising
    \begin{align*}
        q(x_{1 \::\: T} | x_0) &= \prod_{t = 1}^T q(x_t | x_{t - 1}) \\
        q(x_t | x_{t - 1}) &= N(x_t; \sqrt{1 - \beta_t} x_{t - 1}, \beta_t I)
    \end{align*}
    noising schedule: $\beta_t \in [0, 1]$, and follows a predefined schedule
\end{itemize}

\subsection*{DDPM Training}
For each data sample $x_0$ (e.g., image), pick random step $t$
\begin{enumerate}
	\item sample a noise from standard Gaussian $N(0, I) \rightarrow \epsilon$
	\item Noise data at step $t$: $x_t = \sqrt{\bar{\alpha}_t} x_0 + \sqrt{1 - \bar{\alpha}_t} \epsilon, \bar{\alpha}_t = \prod_{s = 1}^t (1 - \beta_s)$
	\item Estimate noise using neural network: $\hat{\epsilon}_\theta = \text{NN}_{\text{forward}}(x_t, t)$
	\item Compute MSE loss: $L = \Vert \epsilon - \hat{\epsilon}_\theta\Vert^2$
	\item Compute gradient via backward through NN
	\item Update model parameters using optimization algorithm (e.g., Adam)
\end{enumerate}
\textbf{DDPM training is essentially predicting the noise based on image}.

\subsection*{DDPM Inference}
\begin{enumerate}
	\item Randomly sample from Gaussian(0, I)
	\item For step $t = T$ to 1:
	\begin{itemize}
	    \item Estimate noise using neural network: $\hat{\epsilon}_\theta = NN_{\text{forward}}(x_t, t)$
	    \item Generate noise from $N(0, I) \rightarrow \eta$
	    \item Estimate data mean: $\tilde{\mu}_{t - 1} = \frac{1}{\sqrt{\alpha_t}}\left(x_t - \frac{\beta_t}{\sqrt{1 - \bar{\alpha}_t}} \hat{\epsilon}_\theta\right)$
	    \item Update data: $x_{t - 1} = \tilde{\mu}_{t - 1} + \sqrt{\beta_t} \eta$
    \end{itemize}
\end{enumerate}
\textbf{DDPM inference is essentially iteratively removing predicted noise.}

\subsection*{Intuition and Theory}
\begin{itemize}
	\item Training objective:
\end{itemize}
\begin{align*}
    L &= -\log p_\theta(x_0) \\
    &\geq -E \left[\log \frac{p_\theta(x_{0 \::\: T})}{q(x_{1 \::\: T} | x_0)}\right] = L_{elbo} \\
    &= \sum_{t = 1}^T \text{KL}(q(x_{t - 1} | x_t, x_0) || p_\theta(x_{t - 1} | x_t)) + \text{KL}(q(x_T | x_0) || p_\theta(x_T)) - E \log p_\theta(x_0 | x_1)
\end{align*}
\begin{itemize}
	\item $L_{elbo}$ is ELBO loss. 
	\item We can use Jensen's inequality, posterior for linear Gaussian distribution
\end{itemize}
\begin{align*}
    q(x_t | x_0) &= N(x_t; \sqrt{\bar{\alpha}_t} x_0, (1 - \bar{\alpha}_t)I) \\
    \bar{\alpha}_t &= \prod_{s = 1}^t (1 - \beta_s) \\
    & \\
    q(x_{t - 1} | x_t, x_0) &= N(x_{t - 1}; \tilde{\mu}_t, \tilde{\beta}_t I) \\
    \tilde{\mu}_t &= \frac{1}{\sqrt{\alpha_t}} \left(x_t - \frac{\beta_t}{\sqrt{1 - \bar{\alpha}_t}} \epsilon\right) \\
    \tilde{\beta}_t &= \frac{1 - \bar{\alpha}_{t - 1}}{1 - \bar{\alpha}_t} \\
    & \\
    \intertext{What about $p_\theta(x_{t - 1} | x_t)$?  Just set:}
    p_\theta(x_{t - 1} | x_t) &= N(\mu_\theta(x_t, t), \sigma_t^2 I)\\\\
\end{align*}
\begin{align*}
    &\text{KL}(q(x_{t - 1} | x_t, x_0) || p_\theta(x_{t - 1} | x_t)) \\
    &= \text{KL}(N(x_{t - 1}; \tilde{\mu}_t, \tilde{\beta}_t I) || N(\mu_\theta(x_t, t), \sigma_t^2 I)) \\
    \intertext{let $\sigma_t^2 = \tilde{\beta}_t$, and let}
    \mu_\theta(x_t, t) &= \frac{1}{\sqrt{\alpha_t}} \left(x_t - \frac{\beta_t}{\sqrt{1 - \bar{\alpha}_t}} \hat{\epsilon}_\theta (x_t, t)\right) \\
\end{align*}
\[\min \Vert \tilde{\mu}_t - \mu_\theta(x_t, t) \Vert^2\]
is equivalent to
\[\min \frac{\beta_t^2}{\alpha_t (1 - \bar{\alpha}_t)} \Vert \epsilon - \hat{\epsilon}_\theta(x_t, t) \Vert^2\]

Essentially, DDPM is using NN to predict noise given noised data

\subsection*{Summary}
Diffusion models:
\begin{itemize}
	\item is a probabilistic generative model by iteratively removing noise
	\item construct a sequence of corrupted data by applying Gaussian noises
	\item use a neural network to estimate the noise given the corrupted data (without knowing the original clean data)
	\item once estimate the noise, use posterior estimation to remove noise
	\item can use any suitable NN (e.g., UNet)
	\item much better than VAE, but much slower
\end{itemize}

\section*{Multi-Objective Drug Molecule Generation via MARS}
Drug discovery and development is a long and risky process.  The successful launch of a new drug cost \$1.3 Billion on average and more than 10 years.

\subsection*{In silico Drug Discovery}
\begin{itemize}
	\item Using computers to predict drug activities.
	\item Faster and cheaper.
	\item But drug candidates are either existing molecules or manually designed ones.
\end{itemize}

\subsubsection*{Discovering Treatment for COVID-19}
\begin{itemize}
	\item Paxlovid
	\item Nirmatrelvir (by modifying prior protease inhibitor)
	\item Ritonavir: originally for HIV
	\item Remdesivir: originally for hepatitis C
\end{itemize}

\subsubsection*{AI Powered Drug Discovery}
\begin{itemize}
	\item AI can generate drug-like molecules with desired properties.
	\item Using deep generative model (GENTRL) to design novel and biological active drug candidates in 46 days.
\end{itemize}
\begin{center} 
	\includegraphics*[width=0.9\textwidth]{W3_5.png} 
\end{center}

\subsubsection*{Generative Models for Molecules}
\begin{center} 
	\includegraphics*[width=0.8\textwidth]{W3_6.png} 
\end{center}

\subsection*{Perspective 1: Molecule Sequence Generation}
\textbf{Transformer based:}
\begin{center} 
	\includegraphics*[width=0.8\textwidth]{W3_7.png} 
\end{center}
\[p(x) = p(x_1) p(x_2 | x_1) p(x_3 | x_1, x_2) \cdots\]
\textbf{Generation from Latent Space:}
\begin{center} 
	\includegraphics*[width=0.9\textwidth]{W3_8.png} 
\end{center}
\textbf{Reinforcement-Learning based:}
\begin{center} 
	\includegraphics*[width=0.9\textwidth]{W3_9.png} 
\end{center}
\textbf{RationaleRL}
\begin{center} 
	\includegraphics*[width=0.7\textwidth]{W3_10.png} 
\end{center}

\subsubsection*{Challenges}
\begin{itemize}
	\item Current generative models can only find molecules with high similarities comparing to existing drugs / potential drug candidates.
\end{itemize}
\begin{center} 
	\includegraphics*[width=0.8\textwidth]{W3_11.png} 
\end{center}

\subsection*{Practical Drug Discovery}
Broader exploration of chemical space
\begin{itemize}
	\item Chemical space is huge but under-explored.
	\item About 85\% drug candidates fail at clinical trials
	\item Finding diverse drug candidates can improve the success rate in clinical trials.
\end{itemize}

\subsection*{Multi-Objective Optimization for Drug Discovery}
\begin{itemize}
	\item Successful drugs need to meet multiple requirements, e.g., bioactivity, safety, and etc.
	\item In drug discovery, these properties should also be optimized.
\end{itemize}
\textbf{Goal: Finding Diverse Molecules Satisfying Multiple Objectives}
\begin{itemize}
	\item Satisfy multiple properties with high scores
	\item Produce diverse and novel molecules
	\item Does not rely on lab measured data
\end{itemize}

\subsection*{MARS: Markov Molecule Sampling}
\begin{itemize}
	\item MCMC sampling from
	\[\pi(x) = s_1(x) \circ s_2(x) \circ s_3(x) \circ \cdots \circ s_K(x)\]
    \begin{itemize}
	    \item Where each $s_1, s_2, \dots$ represent desired properties.
	    \item $s_k(x)$ is a scoring function for a given desired property.
	    \begin{itemize}
	        \item QED: drug-likeness
	        \item SA: possibility to synthesize
        \end{itemize}
        \item $\pi(x)$ is the unnormalized distribution over molecule space.
    \end{itemize}
\end{itemize}
MARS does \textbf{Iterative Graph Editing}.
\begin{center} 
	\includegraphics*[width=0.9\textwidth]{W3_12.png} 
\end{center}
\begin{itemize}
	\item Start from initial molecule (e.g., Ethane)
	\item For each step, propose a new molecule $x'$ by modifying existing one, $x' \sim q(x' | x^{(t - 1)})$
	\item Accept wrt ratio:
	\[\mathcal{A}(x, x') = \min \left\{\frac{\pi^\alpha (x') q(x | x')}{\pi^\alpha(x) q(x' | x)}\right\}\]
\end{itemize}

\subsubsection*{Adaptive Proposal Learning in MARS}
\begin{itemize}
	\item MCMC sampling
	\item Markov-chain annealing scheme to find optimal samples
	\item Adaptive MCMC proposal: MPNN learned on sampled molecular graphs
\end{itemize}

\subsection*{Molecular Graph Editing}
\begin{itemize}
	\item Adding fragment
	\begin{center} 
	    \includegraphics*[width=0.65\textwidth]{W3_13.png} 
    \end{center}
    \item Deleting fragment
    \begin{center} 
	    \includegraphics*[width=0.65\textwidth]{W3_14.png} 
    \end{center}
\end{itemize}

\subsection*{Learning Proposal Network (MPNN)}
Modeling editing actions as node, edge, and graph prediction problems
\begin{center} 
	\includegraphics*[width=0.7\textwidth]{W3_15.png} 
\end{center}

\subsection*{MARS with Docking Score}
Grade MARS molecules by how well it binds to protein targets
\begin{center} 
	\includegraphics*[width=0.5\textwidth]{W3_16.png} 
\end{center}

\subsubsection*{Integrating 3D Structures}
\begin{center} 
	\includegraphics*[width=0.8\textwidth]{W3_17.png} 
\end{center}

\subsubsection*{Adaptive Proposal Training}
\begin{verbatim}
Initialize molecules x_i^{(0)} for i = 1 to N
Initialize molecular editing model M_theta
Create an empty editing model training dataset D = {};
for t = 1, 2, ..., do
    for i = 1, 2, ..., N, do
        calculate probability distributions (p_add, p_frag, p_del) = M_theta(x_i^{t - 1})
        Sample a candidate molecule x' from the proposal distribution
            q(x' | x_i^{t - 1}) defined with probability distributions p_add, p_frag, p_del
        if u < A(x_i^{t - 1}, x') where u ~ U_{[0, 1]} then
            Accept the candidate molecule x_i^t = x'
        else
            Refuse the candidate molecule x_i^t = x^{t - 1}
        end
        if the candidate improves, i.e., pi(x') > pi(x_i^{t - 1}) then
            adding the editing record (x_i^{t - 1}, x') into the dataset D
        end
    end
    Update model M_theta with the current dataset D in a MLE manner
end
\end{verbatim}

\subsubsection*{MARS Generates Better Molecules}
\begin{center} 
	\includegraphics*[width=0.8\textwidth]{W3_18.png} 
\end{center}
\begin{itemize}
	\item Generation Objectives
	\begin{itemize}
	    \item GSK3$\beta$, JNK3: biological inhibition
	    \item QED: drug-likeness score
	    \item SA: synthetic accessibility score
    \end{itemize}
    \item Evaluate with Product score = success\_rate $\times$ novelty $\times$ diversity
\end{itemize}

\subsubsection*{MARS Explores Large Chemical Space}
\begin{center} 
	\includegraphics*[width=0.9\textwidth]{W3_19.png} 
\end{center}

\subsection*{Novel Compounds Found}
\begin{center} 
	\includegraphics*[width=0.8\textwidth]{W3_20.png} 
\end{center}

\subsection*{Key Takeaways}
\begin{itemize}
	\item Challenges in Drug discovery: multi-objective, novel and diverse, lack of data
	\begin{itemize}
	    \item Modeled as a molecule generation problem
	    \begin{itemize}
	        \item Based on MCMC sampling
	        \item Self-adaptive proposal trained on the fly $\Rightarrow$ no need for data
	        \item Generates better molecules and explores larger chemical space
        \end{itemize}
        \item $\Rightarrow$ can discover novel and diverse drug-like molecules
        \item Challenges remaining:
        \begin{itemize}
	        \item More properties
	        \item Larger molecule, peptide, protein
        \end{itemize}
    \end{itemize}
\end{itemize}

\section*{Decoding Single-Cell Spatial Genomics and Epigenomics with AI/ML}
\subsection*{Introduction}
Genome function in different cell types
\begin{itemize}
	\item Numerous different cell types with distinct functions in our body
	\item Same genome but different gene regulation and epigenome
	\item Development of single cell technology provides closer look at cell types and cell states
\end{itemize}
Multiscale cellular structure and function
\begin{itemize}
	\item Single-cell 3D epigenome and gene regulation
	\item Cellular spatial organization and interaction
	\item Graph / hypergraph neural network, self-supervised, latent embedding, metagene, attention, (foundation model)
\end{itemize}
Large-scale genome organization in the nucleus
[FILL 4]

\subsubsection*{Hi-C Contact Map}
[FILL 5]

\subsubsection*{Multiscale 3D genome organization}
[FILL 6]

\subsubsection*{Challenges for Single-cell 3D genome analysis using scHi-C data}
\begin{itemize}
	\item High dimensionality
	\item Sparse and noisy data
	\item Multiscale 3D genome features
\end{itemize}
What do 3D genome look like in \textbf{single cells}, and how do multiscale 3D features vary in individual cells and different cell types?

\subsection*{Hypergraphs for Higher-Order Interactions}
\begin{itemize}
	\item Hypergraphs are used to represent higher-order interactions
	\begin{itemize}
	    \item Example: events (human, location, activity)
    \end{itemize}
    \item A \textbf{hypergraph} $H = (V, E)$
    \begin{itemize}
	    \item $V / E$: the set of nodes / \textbf{hyperedges}
	    \item $e \in E$: a hyperedge connects two or more nodes
	    \item $\forall e \in E, |e| = k$, $H$ is a $k$-uniform hypergraph
    \end{itemize}
    \item \textbf{Hyper-SAGNN} - hypergraph representation learning
    \begin{itemize}
	    \item Input: a hypergraph with features for each node (non-uniform heterogeneous hypergraph)
        \item We aim to:
        \begin{itemize}
	        \item Learn embeddings for the nodes in the hypergraph
	        \item Learn to predict the existence of hyperedges given the node embeddings
        \end{itemize}
    \end{itemize}
\end{itemize}
[FILL 8, rearrange]

\subsubsection*{Higashi - modeling scHi-C data as a hypergraph}
[FILL 9]
[FILL 10]

\textbf{Higashi separates complex cell types in the human prefrontal cortex}
\begin{itemize}
	\item Higashi embeddings separate neuron subtypes using only scHi-C part of sn-m3c-seq (data from Lee et al. \textit{Nat Methods} 2019)
	\item Cell type-specific 3D chromatin structures near marker genes
\end{itemize}
[FILL 11, rearrange]

\subsection*{Mechanisms of Genome Folding}
[FILL 12, all]

\subsection*{Patterns of Chromosome Spatial Segregation based on Hi-C data}
\begin{itemize}
	\item Chromosomes are segregated into A and B compartments
	\item High-coverage Hi-C in GM12878 identified \textbf{subcompartments} by clustering inter-chromosomal contact maps
\end{itemize}
[FILL 13]

\subsection*{SNIPER - Inferring Hi-C Subcompartments}
\begin{itemize}
	\item Autoencoder compresses sparse inter-chromosomal Hi-C contacts into embeddings and imputes missing contacts
	\item A classifier then uses embedded interchromosomal Hi-C data to predict subcompartments
\end{itemize}

\subsection*{scGHOST: Identifying Single-Cell 3D Genome Subcompartments}
\begin{itemize}
	\item scHi-C contact maps as \textbf{graphs}
	\item Cell embeddings define k-NN cells.  \textbf{Genomic locus embeddings} lead to subcompartments
	\item A unique \textbf{random sampling} procedure that filters noise in imputed scHi-C data
\end{itemize}
[FILL 15]
[FILL 16]
\begin{itemize}
	\item Sampling based on both first-order random walks and second-order random walks
	\item Graph node embedding using NN
\end{itemize}

\subsubsection*{Single-Cell Subcompartments of Human Prefrontal Cortex}
[FILL 17]

\subsubsection*{Spatial Organization of Genomes and Cells}
[FILL 18]

\subsubsection*{Spatial Transcriptomic Technologies}
Imaging-based - single cell resolution
\begin{itemize}
	\item STARmap / STARmap PLUS (Wang et al. \textit{Science} 2018 / Zeng et al. \textit{Nat Neurosci} 2023)
	\item seqFISH+ (Eng et al. \textit{Nature} 2019)
	\item MERFISH / Vizgen MERSCOPE (Moffitt et al. \textit{Science} 2018)
\end{itemize}
Sequencing-based - full transcriptome
\begin{itemize}
	\item Visium (Stahl et al. \textit{Science} 2016)
	\item Slide-seq / Slide-seq V2 (Rodriques et al. \textit{Science} 2019 / Stickels et al. \textit{Nat Biotech} 2021)
	\item Stereo-seq (Chen et al. \textit{Cell} 2022)
\end{itemize}

\end{document}